\originalchapter{红巨星演化、氦闪与中微子习题}
本章对应8.901 PS5.2--4。原题给的经验函数与理想热模型均实际积分，解答为独立补写。它们是研究某个反馈或尺度的练习，不是完整恒星演化代码；额外参数和原式问题分别说明。

\src{8.901 PS5.2，红巨星支上的演化。}
\begin{exercise}
低质量恒星离开主序后，薄氢壳燃烧围绕简并氦核。题给$L=2\times10^5\Ls(m_c/\Ms)^6$、$R=3700\Rs(m_c/\Ms)^4$。用$L=0.007\dot m_cc^2$，从$m_c=0.1\Ms$积分到$0.45\Ms$，画$R,L,T_{\rm eff}$随时间，以及$\log L$对$\log T_{\rm eff}$的HR轨迹，温度向左增加。
\end{exercise}
\begin{solution}
令$x=m_c/\Ms$，$C=2\times10^5\Ls/(0.007\Ms c^2)$，则$\dot x=Cx^6$。积分给
\[
 x(t)=[x_i^{-5}-5Ct]^{-1/5},\quad
 t_f=\frac{x_i^{-5}-x_f^{-5}}{5C}\simeq1.035\times10^{10}\ \mathrm{yr}.
\]
把$x(t)$代回得到$L,R$；$T_{\rm eff}=T_\odot(L/\Ls)^{1/4}(R/\Rs)^{-1/2}\propto x^{-1/2}$。核增长时整星光度和半径迅速增加，表面反而冷却。取$T_\odot=5772$ K，初点$(L/\Ls,R/\Rs,T)=(0.2,0.37,6346)$，终点约$(1661,151.7,2991)$。
\begin{center}
\begin{tikzpicture}\begin{axis}[width=.47\textwidth,height=4.5cm,xlabel={$t$（yr）},ylabel={太阳单位},ymode=log,legend pos=north west]
\addplot[structurecolor] table[x=t_yr,y=L_solar,col sep=comma]{qa/stellar-giant-evolution.csv};\addlegendentry{$L/\Ls$}
\addplot[black,dashed] table[x=t_yr,y=R_solar,col sep=comma]{qa/stellar-giant-evolution.csv};\addlegendentry{$R/\Rs$}
\end{axis}\end{tikzpicture}\hfill
\begin{tikzpicture}\begin{axis}[width=.47\textwidth,height=4.5cm,xlabel={$t$（yr）},ylabel={$T_{\rm eff}$（K）}]
\addplot[structurecolor] table[x=t_yr,y=Teff_K,col sep=comma]{qa/stellar-giant-evolution.csv};
\end{axis}\end{tikzpicture}
\end{center}
\begin{center}\begin{tikzpicture}\begin{axis}[width=.55\textwidth,height=4.4cm,xmode=log,ymode=log,x dir=reverse,xlabel={$T_{\rm eff}$（K，左侧更热）},ylabel={$L/\Ls$}]
\addplot[structurecolor,thick] table[x=Teff_K,y=L_solar,col sep=comma]{qa/stellar-giant-evolution.csv};
\end{axis}\end{tikzpicture}\end{center}
极长时间主要来自$x=0.1$附近极低$L$的外推，而该初点连半径都不足太阳，不能说它是真实已膨胀红巨星。答案遵照题给函数，同时指出其低核质量端的适用性问题；不通过改起点缩短时间而仍声称解了原题。
\end{solution}

\src{8.901 PS5.3，氦闪，原题全部条件与数值要求。}
\begin{exercise}
一克纯氦初始$\rho=2\times10^5\ \mathrm{g\,cm^{-3}}$、$T_0=1.5\times10^8$ K。只考虑三氦反应，保持密度、氦质量分数$Y=1$且无热流，题给
\[
 \epsilon_{3\alpha}=5.1\times10^8\rho^2Y^3T_9^{-3}\ee^{-4.4027/T_9}
 \quad\mathrm{erg\,g^{-1}s^{-1}},\quad T_9=T/10^9\ \mathrm K.
\]
电子比热$C_{V,e}=1.35\times10^5T x\sqrt{1+x^2}/\rho$，$\rho=2\times10^6x^3$；离子比热按原式$C_{V,\mathrm{He}}=3\kb/(4m_p)$。实际求$T(t)$，到非相对论简并分界$\rho=10^{-8}T^{3/2}$停止，画时间以天为单位的温度曲线；给氦闪时间、数值精度检查，并说明结果的物理合理性。
\end{exercise}
\begin{solution}
固定$\rho$使$x=(0.1)^{1/3}=0.46416$。热方程$C_V(T)\dot T=\epsilon(T)$，其中$C_V=C_{V,e}+C_{V,\mathrm{He}}$。温度单调增加，可将时间微分方程改写成
\[
 t(T)=\int_{T_0}^T\frac{C_V(T')}{\epsilon(T')}\dd T'.
\]
这仍是实际数值计算，且无需用固定时间步跨过升温末段。分界$T_d=(\rho/10^{-8})^{2/3}=7.368\times10^8$ K。自适应求积相对误差$10^{-10}$所得$t(T_d)=667284.2$ s，即7.7232 d；温度加倍时667265.3 s，距分界约19 s。若定义“闪开始”为温度加倍，则为7.7230 d；模型没有天然不连续的开始时刻，必须给定义。
\begin{center}\begin{tikzpicture}\begin{axis}[width=.72\textwidth,height=4.8cm,xlabel={$t$（d）},ylabel={$T$（K）},legend pos=north west]
\addplot[structurecolor] table[x expr=\thisrow{t_printed_s}/86400,y=T_K,col sep=comma]{qa/stellar-helium-flash.csv};\addlegendentry{原题比热}
\addplot[black,dashed] table[x expr=\thisrow{t_ideal_ion_s}/86400,y=T_K,col sep=comma]{qa/stellar-helium-flash.csv};\addlegendentry{单原子离子比热修正}
\end{axis}\end{tikzpicture}\end{center}
纯氦离子的经典定容热容为$3\kb/2$，每克离子数$1/(4m_p)$，所以理想离子比热应为$3\kb/(8m_p)$，原式大两倍。不能静默修正后只给一个时间。另算修正式得到$t(T_d)=489384.3$ s，即5.6642 d。图中两条曲线分别保留题给计算和物理比热判断。

升温增强反应率，而简并电子压力暂时不强依赖温度，缺少普通热膨胀降温反馈，因而可有热失控。固定密度、无热流、固定$Y$在真实氦闪中会失效；停止边界也只是近似简并判据。应重算膨胀、燃料和输运，不能把这几天理解为对某颗红巨星实际爆发日期的预测。
\end{solution}

\src{8.901 PS5.4，中微子振荡，(a)--(b)。}
\begin{exercise}
两种味态$|\nu_e\rangle=\cos\theta|\nu_1\rangle+\sin\theta|\nu_2\rangle$、$|\nu_\mu\rangle=-\sin\theta|\nu_1\rangle+\cos\theta|\nu_2\rangle$，质量态能量$E_j=\sqrt{p^2c^2+m_j^2c^4}$。$t=0$在原点产生动量$p$的纯电子中微子，向$z$传播。
(a)用自由平面波、小质量展开求波函数，提出共同相位，取$E\simeq pc,z\simeq ct$；
(b)证明转成缪子味的概率$P=\sin^2(2\theta)\sin^2[c^3\Delta m^2z/(4\hbar E)]$和振荡长度。原作业提示太阳内部MSW效应，引用2004年$\sin^22\theta=0.82\pm0.07$、$\Delta m^2c^4=(8.2\pm0.6)\times10^{-5}\ \mathrm{eV^2}$。
\end{exercise}
\begin{solution}
小质量时$E_j\simeq pc+m_j^2c^3/(2p)$。共同平面相位$\exp[-\ii pc(t-z/c)/\hbar]$不影响味测量，移除后
\[
 |\psi(z)\rangle=\ee^{-\ii m_1^2c^3z/(2\hbar E)}\cos\theta|\nu_1\rangle
 +\ee^{-\ii m_2^2c^3z/(2\hbar E)}\sin\theta|\nu_2\rangle.
\]
投影到$\nu_\mu$得振幅$\sin\theta\cos\theta(\ee^{-\ii\phi_2}-\ee^{-\ii\phi_1})$；用$|\ee^{\ii a}-\ee^{\ii b}|^2=4\sin^2[(a-b)/2]$即得题式。概率周期使正弦自变量增加$\pi$，所以
\[
 L_{\rm osc}=\frac{4\pi\hbar E}{c^3|\Delta m^2|}.
\]
需要非零混合与质量差；绝对质量共同增量只改共同相位，振荡不单独确定最轻态质量。

太阳电子背景对电子味多一个前向散射势，改变传播本征态和有效混合；密度缓变可有MSW增强。因此太阳缺失量不必在所有能量都只由地日几何相位解释。两味平面波近似忽略波包分离、三味和介质非均匀性；2004年数值保留年代，不作当前最佳值。
\end{solution}
